\documentclass[aspectratio=169]{beamer}

\usepackage{amsmath}
\usepackage{amssymb}
\usepackage{bm}
\usepackage{mathtools}

\title{Terminal-State Trajectory Formulation for Dynamic Throwing}
\subtitle{Ballistic planning, finite-time motion generation, and dual-arm control}
\author{IITGN ROBOTICS}
\institute{IIT Gandhinagar}

\newcommand{\vect}[1]{\bm{#1}}
\newcommand{\R}{\mathbb{R}}
\newcommand{\diag}{\operatorname{diag}}

\begin{document}

%------------------------------------------------
\begin{frame}
    \titlepage
\end{frame}

%------------------------------------------------
\begin{frame}{Goal of the New Formulation}

The object must pass through a prescribed release state,
not settle at a stationary equilibrium:
\[
\boxed{
\vect{x}(T)=\vect{x}_{\mathrm{rel}},
\qquad
\dot{\vect{x}}(T)=\vect{v}_{\mathrm{rel}}
}
\]

\begin{itemize}
    \item $\vect{x}_{\mathrm{rel}}$: desired object position at release.
    \item $\vect{v}_{\mathrm{rel}}$: velocity required to reach the landing point.
    \item $T$: duration of the controlled pre-release motion.
    \item At $t=T$, the object is moving; it is not at rest.
\end{itemize}

The formulation therefore has two parts:
\begin{enumerate}
    \item Compute the ballistic terminal velocity $\vect{v}_{\mathrm{rel}}$.
    \item Construct a finite-time trajectory that reaches
    $(\vect{x}_{\mathrm{rel}},\vect{v}_{\mathrm{rel}})$.
\end{enumerate}

\end{frame}

%------------------------------------------------
\begin{frame}{Why the Earlier Autonomous DS Is Inappropriate}

The earlier formulation had the form
\[
\ddot{\vect{x}}
=-\vect{K}_{\mathrm{DS}}(\vect{x}-\vect{x}_{\mathrm{rel}})
-\vect{B}_{\mathrm{DS}}(\dot{\vect{x}}-\vect{v}_{\mathrm{rel}}).
\]

At an equilibrium, $\dot{\vect{x}}=\vect{0}$ and
$\ddot{\vect{x}}=\vect{0}$. Therefore,
\[
\vect{K}_{\mathrm{DS}}(\vect{x}_{\mathrm{eq}}-\vect{x}_{\mathrm{rel}})
=\vect{B}_{\mathrm{DS}}\vect{v}_{\mathrm{rel}},
\]
and hence
\[
\boxed{
\vect{x}_{\mathrm{eq}}
=\vect{x}_{\mathrm{rel}}
+\vect{K}_{\mathrm{DS}}^{-1}\vect{B}_{\mathrm{DS}}
\vect{v}_{\mathrm{rel}},
\qquad
\vect{v}_{\mathrm{eq}}=\vect{0}
}
\]

Thus, for $\vect{v}_{\mathrm{rel}}\neq\vect{0}$, the desired release
state is not an equilibrium of the earlier DS. A terminal trajectory is
needed instead.

\end{frame}

%------------------------------------------------
\begin{frame}{Step 1: Define the Ballistic Geometry}

Let
\[
\vect{p}_{\mathrm{rel}}
=\begin{bmatrix}x_{\mathrm{rel}}&y_{\mathrm{rel}}&z_{\mathrm{rel}}\end{bmatrix}^{T},
\qquad
\vect{p}_{\mathrm{land}}
=\begin{bmatrix}x_{\mathrm{land}}&y_{\mathrm{land}}&z_{\mathrm{land}}\end{bmatrix}^{T}.
\]

The displacement is
\[
\Delta\vect{p}=\vect{p}_{\mathrm{land}}-\vect{p}_{\mathrm{rel}}.
\]

Its horizontal range and horizontal direction are
\[
R=\left\|\Delta\vect{p}_{xy}\right\|,
\qquad
\hat{\vect{e}}_{h}
=\frac{\Delta\vect{p}_{xy}}{R},
\]
where $\Delta\vect{p}_{xy}=[\Delta p_x,\Delta p_y]^T$.

The vertical displacement is
\[
\Delta z=z_{\mathrm{land}}-z_{\mathrm{rel}}.
\]

The world-up direction is fixed as
\[
\hat{\vect{z}}=\begin{bmatrix}0&0&1\end{bmatrix}^{T}.
\]

\end{frame}

%------------------------------------------------
\begin{frame}{Step 2: Projectile Motion Under Gravity}

During free flight, gravity is
\[
\vect{g}=\begin{bmatrix}0&0&-g\end{bmatrix}^{T},
\qquad g>0.
\]

Ignoring aerodynamic drag, the object position is
\[
\boxed{
\vect{p}(t)
=\vect{p}_{\mathrm{rel}}
+\vect{v}_{\mathrm{rel}}t
+\frac{1}{2}\vect{g}t^2
}
\]

For a chosen launch angle $\theta$, write the release velocity as
\[
\boxed{
\vect{v}_{\mathrm{rel}}
=v_{\mathrm{rel}}\cos\theta
\begin{bmatrix}\hat{\vect{e}}_{h}\\0\end{bmatrix}
+v_{\mathrm{rel}}\sin\theta\,\hat{\vect{z}}
}
\]

Here $v_{\mathrm{rel}}=\|\vect{v}_{\mathrm{rel}}\|$ is the unknown
release speed.

\end{frame}

%------------------------------------------------
\begin{frame}{Step 3: Derive the Required Release Speed}

Horizontal motion gives the flight time
\[
R=v_{\mathrm{rel}}\cos\theta\,t_f
\quad\Longrightarrow\quad
t_f=\frac{R}{v_{\mathrm{rel}}\cos\theta}.
\]

Vertical motion at landing gives
\[
\Delta z
=v_{\mathrm{rel}}\sin\theta\,t_f
-\frac{1}{2}gt_f^2.
\]

Substituting the horizontal flight time,
\[
\Delta z
=R\tan\theta
-\frac{gR^2}{2v_{\mathrm{rel}}^2\cos^2\theta}.
\]

Solving for the release speed,
\[
\boxed{
v_{\mathrm{rel}}
=\sqrt{
\frac{gR^2}
{2\cos^2\theta\left(R\tan\theta-\Delta z\right)}
}
}
\]

A real solution requires
$R\tan\theta-\Delta z>0$.

\end{frame}

%------------------------------------------------
\begin{frame}{Step 4: Measure the Initial Throw State}

When the controller enters the throw phase at time $t_0$, it measures
the actual object state:
\[
\boxed{
\vect{x}_0=\vect{x}(t_0),
\qquad
\vect{v}_0=\dot{\vect{x}}(t_0)
}
\]

The desired terminal state is
\[
\boxed{
\vect{x}_f=\vect{x}_{\mathrm{rel}},
\qquad
\vect{v}_f=\vect{v}_{\mathrm{rel}}
}
\]

The pre-release trajectory must satisfy four boundary conditions:
\[
\vect{x}_{\mathrm{ref}}(0)=\vect{x}_0,
\qquad
\dot{\vect{x}}_{\mathrm{ref}}(0)=\vect{v}_0,
\]
\[
\vect{x}_{\mathrm{ref}}(T)=\vect{x}_{\mathrm{rel}},
\qquad
\dot{\vect{x}}_{\mathrm{ref}}(T)=\vect{v}_{\mathrm{rel}}.
\]

Using the measured state prevents a discontinuity at throw-phase entry.

\end{frame}

%------------------------------------------------
\begin{frame}{Step 5: Normalize the Trajectory Time}

Define elapsed and normalized time as
\[
\tau=t-t_0,
\qquad
s=\frac{\tau}{T},
\qquad 0\leq s\leq 1.
\]

A cubic polynomial is sufficient because each Cartesian coordinate has
four boundary conditions:
\[
x_{\mathrm{ref}}(s)=a_0+a_1s+a_2s^2+a_3s^3.
\]

The same construction is applied independently to the $x$, $y$, and $z$
components. In vector form, this produces a three-dimensional object
trajectory.

\begin{block}{Why a cubic?}
A cubic contains four coefficients, exactly enough to impose the initial
position, initial velocity, terminal position, and terminal velocity.
\end{block}

\end{frame}

%------------------------------------------------
\begin{frame}{Step 6: Cubic Hermite Position Reference}

Define the Hermite basis functions
\[
\begin{aligned}
h_{00}(s)&=2s^3-3s^2+1, &
h_{10}(s)&=s^3-2s^2+s,\\
h_{01}(s)&=-2s^3+3s^2, &
h_{11}(s)&=s^3-s^2.
\end{aligned}
\]

The reference position is
\[
\boxed{
\vect{x}_{\mathrm{ref}}(s)
=h_{00}(s)\vect{x}_0
+h_{10}(s)T\vect{v}_0
+h_{01}(s)\vect{x}_{\mathrm{rel}}
+h_{11}(s)T\vect{v}_{\mathrm{rel}}
}
\]

The factors of $T$ convert endpoint velocities into position units.
The basis functions are selected so that the four boundary conditions
are satisfied exactly.

\end{frame}

%------------------------------------------------
\begin{frame}{Step 7: Reference Velocity}

Because $s=\tau/T$, differentiation uses
\[
\frac{d}{dt}=\frac{1}{T}\frac{d}{ds}.
\]

Therefore,
\[
\boxed{
\begin{aligned}
\vect{v}_{\mathrm{ref}}(s)
={}&\frac{6s^2-6s}{T}\vect{x}_0
+(3s^2-4s+1)\vect{v}_0\\
&+\frac{-6s^2+6s}{T}\vect{x}_{\mathrm{rel}}
+(3s^2-2s)\vect{v}_{\mathrm{rel}}.
\end{aligned}
}
\]

Checking the endpoints gives
\[
\vect{v}_{\mathrm{ref}}(0)=\vect{v}_0,
\qquad
\vect{v}_{\mathrm{ref}}(1)=\vect{v}_{\mathrm{rel}}.
\]

Hence the object reaches the release point with a nonzero prescribed
velocity instead of stopping there.

\end{frame}

%------------------------------------------------
\begin{frame}{Step 8: Closed-Loop Terminal Tracking}

The polynomial is the feedforward motion reference. Feedback corrects
tracking errors caused by robot dynamics, contacts, gravity, and limits:
\[
\vect{e}_x=\vect{x}_{\mathrm{ref}}-\vect{x},
\qquad
\vect{e}_v=\vect{v}_{\mathrm{ref}}-\vect{v}.
\]

The commanded object velocity is
\[
\boxed{
\vect{v}_{o,\mathrm{cmd}}
=\vect{v}_{\mathrm{ref}}
+\vect{K}_x\vect{e}_x
+\vect{K}_v\vect{e}_v
}
\]

\begin{itemize}
    \item $\vect{v}_{\mathrm{ref}}$: follows the planned motion.
    \item $\vect{K}_x\vect{e}_x$: corrects position error.
    \item $\vect{K}_v\vect{e}_v$: corrects velocity error.
\end{itemize}

For safety, the magnitude may be limited:
\[
\vect{v}_{o,\mathrm{cmd}}
\leftarrow
\min\!\left(1,\frac{v_{\max}}{\|\vect{v}_{o,\mathrm{cmd}}\|}\right)
\vect{v}_{o,\mathrm{cmd}}.
\]

\end{frame}

%------------------------------------------------
\begin{frame}{Step 9: Object Motion to Contact Motion}

Let the commanded object twist be
\[
\vect{\mathcal{V}}_{o,\mathrm{cmd}}
=\begin{bmatrix}\vect{v}_{o,\mathrm{cmd}}\\\vect{\omega}_{o,\mathrm{cmd}}\end{bmatrix}.
\]

For contact $i$, let $\vect{r}_i$ point from the object center to the
contact point. Rigid-body kinematics gives
\[
\boxed{
\vect{v}_{c_i}
=\vect{v}_{o,\mathrm{cmd}}
+\vect{\omega}_{o,\mathrm{cmd}}\times\vect{r}_i
}
\]

Equivalently,
\[
\vect{\mathcal{V}}_{c_i}
=\underbrace{
\begin{bmatrix}
\vect{I}_3&-[\vect{r}_i]_{\times}\\
\vect{0}&\vect{I}_3
\end{bmatrix}}_{\vect{H}_i}
\vect{\mathcal{V}}_{o,\mathrm{cmd}}.
\]

Both contacts receive the complete object motion. A velocity constraint is
not divided by the number of arms. In the translational implementation,
$\vect{\omega}_{o,\mathrm{cmd}}=\vect{0}$.

\end{frame}

%------------------------------------------------
\begin{frame}{Step 10: Add Contact-Level Admittance}

Each hand regulates its grasp force using the measured contact wrench:
\[
\boxed{
\Delta\vect{\mathcal{V}}_{a,i}
=\vect{K}_{a}^{-1}
\left(\vect{F}^{\star}_{i}-\vect{F}_{i}\right)
}
\]

For grasp-force regulation, the translational correction is projected onto
the pad-normal direction $\vect{n}_i$:
\[
\Delta\vect{v}_{a,i}
=\vect{n}_i\vect{n}_i^T
\left[\vect{K}_{a}^{-1}
(\vect{F}^{\star}_{i}-\vect{F}_{i})\right]_{\mathrm{lin}}.
\]

The desired hand twists are then
\[
\boxed{
\begin{bmatrix}
\vect{\mathcal{V}}_{L}^{\star}\\
\vect{\mathcal{V}}_{R}^{\star}
\end{bmatrix}
=
\begin{bmatrix}
\vect{H}_{L}\\
\vect{H}_{R}
\end{bmatrix}
\vect{\mathcal{V}}_{o,\mathrm{cmd}}
+
\begin{bmatrix}
\Delta\vect{\mathcal{V}}_{a,L}\\
\Delta\vect{\mathcal{V}}_{a,R}
\end{bmatrix}
}
\]

Thus, object transport and grasp-force correction are combined additively.

\end{frame}

%------------------------------------------------
\begin{frame}{Step 11: Convert Hand Velocities to Joint Velocities}

Stack the two arm Jacobians:
\[
\vect{J}_{H}
=
\begin{bmatrix}
\vect{J}_{L}&\vect{0}\\
\vect{0}&\vect{J}_{R}
\end{bmatrix},
\qquad
\begin{bmatrix}
\vect{\mathcal{V}}_{L}\\
\vect{\mathcal{V}}_{R}
\end{bmatrix}
=\vect{J}_{H}\dot{\vect{q}}.
\]

The desired joint velocity is computed using a damped least-squares
pseudoinverse:
\[
\boxed{
\dot{\vect{q}}_{\mathrm{des}}
=\vect{J}_{H}^{\#}
\begin{bmatrix}
\vect{\mathcal{V}}_{L}^{\star}\\
\vect{\mathcal{V}}_{R}^{\star}
\end{bmatrix}
}
\]
with
\[
\boxed{
\vect{J}_{H}^{\#}
=\vect{J}_{H}^{T}
\left(\vect{J}_{H}\vect{J}_{H}^{T}
+\lambda^2\vect{I}\right)^{-1}
}
\]

The damping $\lambda>0$ improves numerical robustness near singularities.

\end{frame}

%------------------------------------------------
\begin{frame}{Step 12: Joint-Velocity Tracking}

Define the joint-velocity error
\[
\vect{e}_{\dot q}
=\dot{\vect{q}}_{\mathrm{des}}-\dot{\vect{q}}.
\]

A velocity PID controller produces the commanded joint torque:
\[
\boxed{
\vect{\tau}_{\mathrm{cmd}}
=\vect{K}_{d}\vect{e}_{\dot q}
+\vect{K}_{i}\int\vect{e}_{\dot q}\,dt
+\vect{\tau}_{\mathrm{fric}}
+\vect{\tau}_{\mathrm{bias}}
}
\]

where
\begin{itemize}
    \item $\vect{\tau}_{\mathrm{fric}}$ is friction feedforward,
    \item $\vect{\tau}_{\mathrm{bias}}$ compensates gravity, Coriolis,
          and centrifugal effects,
    \item torque and joint-velocity limits are enforced for safety.
\end{itemize}

Gravity remains active in the physical simulation; bias compensation helps
the robots track the commanded motion while carrying the object.

\end{frame}

%------------------------------------------------
\begin{frame}{Step 13: Terminal Release Condition}

At and after the nominal terminal time $T$, calculate
\[
e_p=\left\|\vect{x}-\vect{x}_{\mathrm{rel}}\right\|,
\qquad
e_v=\left\|\vect{v}-\vect{v}_{\mathrm{rel}}\right\|.
\]

Release is allowed only if
\[
\boxed{
\tau\geq T,
\qquad
e_p<\varepsilon_p,
\qquad
e_v<\varepsilon_v
}
\]

\begin{itemize}
    \item The time condition ensures the terminal reference has been reached.
    \item The position condition ensures release occurs at the planned point.
    \item The velocity condition ensures the required launch velocity has
          actually been achieved.
\end{itemize}

If these conditions are not satisfied, the controller keeps tracking the
terminal state rather than releasing at a poor state.

\end{frame}

%------------------------------------------------
\begin{frame}{Step 14: Opening and Follow-Through}

When release is triggered, grasp-force closure is disabled and the hands
open along opposing outward normals.

For each hand,
\[
\boxed{
\vect{v}_{ee,i}^{\mathrm{release}}
=\vect{v}_{\mathrm{rel}}
+v_{\mathrm{out}}\vect{n}_{\mathrm{out},i}
}
\]

\begin{itemize}
    \item $\vect{v}_{\mathrm{rel}}$ keeps the hand moving with the object.
    \item $v_{\mathrm{out}}\vect{n}_{\mathrm{out},i}$ separates the fingers
          from the object.
    \item Co-motion prevents the hands from abruptly braking the object while
          contact is disappearing.
\end{itemize}

After the short opening interval, the arms enter follow-through while the
object continues independently.

\end{frame}

%------------------------------------------------
\begin{frame}{Step 15: Free Flight After Release}

Once contact is lost, the robot controller no longer determines the object
trajectory. The simulator integrates the physical dynamics.

Under the ideal ballistic model,
\[
\boxed{
\ddot{\vect{p}}=\vect{g}
}
\]
with the release initial conditions
\[
\vect{p}(0)=\vect{p}_{\mathrm{rel}},
\qquad
\dot{\vect{p}}(0)=\vect{v}_{\mathrm{rel}}.
\]

Therefore,
\[
\boxed{
\vect{p}(t)
=\vect{p}_{\mathrm{rel}}
+\vect{v}_{\mathrm{rel}}t
+\frac{1}{2}\vect{g}t^2
}
\]

The controlled phase establishes the release state; gravity then produces
the free-flight trajectory.

\end{frame}

%------------------------------------------------
\begin{frame}{Complete New Formulation}

\begin{enumerate}
    \item Select $\vect{p}_{\mathrm{rel}}$, $\vect{p}_{\mathrm{land}}$,
          and launch angle $\theta$.
    \item Compute $\vect{v}_{\mathrm{rel}}$ from projectile motion.
    \item Measure the actual initial state $(\vect{x}_0,\vect{v}_0)$.
    \item Construct a cubic Hermite trajectory over duration $T$.
    \item Add position and velocity feedback to obtain
          $\vect{v}_{o,\mathrm{cmd}}$.
    \item Map the object twist to both rigid contacts and add grasp
          admittance corrections.
    \item Use the stacked damped Jacobian inverse to obtain
          $\dot{\vect{q}}_{\mathrm{des}}$.
    \item Track the desired joint velocity with bias-compensated torque control.
    \item Release only when both terminal position and velocity are achieved.
    \item Co-move while opening; then allow gravity-driven free flight.
\end{enumerate}

\begin{block}{Central idea}
The release state is a finite-time boundary condition through which the
object passes, not a stationary equilibrium to which it converges.
\end{block}

\end{frame}

%------------------------------------------------
\end{document}
